% Draft: §2, Two analyses and a criterion, for the restructured paper (rewrite
% method agreed 17 Oct 2026; see the evidence repo's DECISIONS.md). Drafted from
% the old draft's §§1, 3.1 and 10, STATE-OF-THE-ARGUMENT (17 Oct), CGEL
% pp. 217, 331, 410--412, 415, 417, 418, 420--421, 434, and Payne et al. (2007)
% pp. 566 n. 2, 569--571 (with n. 4), 582, 588, 600, all checked against the PDFs
% on 17 Oct. Revised after the muse review (reviews/muse-sec2-20261017-183900/).
% Ordinary Head is defined as in STATE: CGEL's tree minus the dependent function
% (the DP or AdjP heads Nom). The old draft's NP projection for determinatives
% (Det:{NP,PP}) is dropped from this paper (Brett, 17 Oct); it belongs to the
% taxonomic paper.
% Bib: NOW and iWeb have no entries yet (L2).
% 18 Oct: footnote on partly/largely added (Matthews 2014 pp. 28--29, checked
% against the text conversion). Constructed examples, judgments Brett's, 18 Oct
% (NEEDS-FROM-BRETT M4, M5): I prefer pure cotton shirts to partly nylon, I
% prefer cotton to partly nylon, This shirt is partly nylon, all good.
% 18 Oct (Brett: "add it"): one sentence after "so both predict that" on why
% a silent head doesn't predict the countability facts (evidence repo
% DECISIONS, comparison with Aarts's empty heads).
% 18 Oct: the p. 217 illustration replaced (Brett: the truncated quotation was
% misleading; CGEL's sentence ends "the subject of quite competent", and Ed is
% the subject of the passive clause). Now the active She considered Ed quite
% competent, with CGEL's term predicand, rechecked against the PDF. Predicand
% glossed as semantic (Brett, 18 Oct): CGEL p. 252, "This semantic predicate
% applies to the predicand, generally subject or object".
% "This lava is partly magma": supplied by Brett, 18 Oct, from
% https://www.conserve-energy-future.com/geothermalenergy.php (no date shown).
%
% Labels it relies on: sec:det-head (§3), sec:mod-head (§4), sec:once (§5),
% sec:fused-relatives (§7). Figure needs forest + the main preamble's \synnode
% and `nominal tree' style. Not yet \input into the main file.

\section{Two analyses and a criterion}\label{sec:analyses}

The framework keeps categories apart from functions. A category is a set of words or phrases, such as noun, determinative, or adjective. A function is the relation a constituent bears within its phrase, such as head, determiner, or modifier. The two vary independently: \mention{nylon} is a noun functioning as head in \mention{made of \bk{nylon}} and as modifier in \mention{a \bk{nylon} shirt} (\textit{CGEL}, p.~421), and an NP functions as determiner in \mention{\bk{Kim's} book}.

So \term{determinative} names a category (\mention{the}, \mention{this}, \mention{some}, \mention{every}, \mention{many}), and \term{determiner} names a function. Figure~\ref{fig:analyses}a gives \textit{CGEL}'s structure for \mention{some apples}. The determinative heads a determinative phrase (DP) that functions as determiner (Det) of the NP.\footnote{DP here is \textit{CGEL}'s determinative phrase, a phrase whose head is a determinative, not the DP of the DP hypothesis, on which the determinative heads the whole noun phrase. \textcite{millerpullum2022tete} draw the same distinction. PHP (p.~566 n.~2), like \textit{CGEL}, take the noun as head of the NP.} The noun heads the \term{nominal} (Nom), the unit between noun and NP, and the nominal is the NP's head. \textit{CGEL} recognizes the unit because \mention{old man} in \mention{the old man} is neither a noun nor an NP, since it can't function as subject or object (p.~329). I write \enquote{nominal} in prose and Nom in trees and quotations. The NP's \term{ultimate head} is the word reached by following head relations down from the NP: \mention{apples} in Figure~\ref{fig:analyses}a, \mention{some} in b and~c.

\begin{figure}[H]
\centering
\begin{minipage}[t]{.36\linewidth}\centering
(a) \mention{some apples}\par\smallskip
\begin{forest} nominal tree
[NP
 [{\synnode{Det}{DP}}
  [{\synnode{Head}{D}}, head edge [\mention{some}]]]
 [{\synnode{Head}{Nom}}, head edge
  [{\synnode{Head}{N}}, head edge [\mention{apples}]]]]
\end{forest}
\end{minipage}\hfill
\begin{minipage}[t]{.3\linewidth}\centering
(b) \mention{some}, fusion\par\smallskip
\begin{forest} nominal tree
[NP
 [{\synnode{Head}{Nom}}, head edge, before drawing tree={x+=1.5em}
  [{\synnode{Det--Head}{DP}}, no edge
   [{\synnode{Head}{D}}, head edge [\mention{some}]]]
  {\draw[-] (!uu.south) -- (); \draw[-,line width=1pt] (!u.south) -- (); }]]
\end{forest}
\end{minipage}\hfill
\begin{minipage}[t]{.3\linewidth}\centering
(c) \mention{some}, ordinary Head\par\smallskip
\begin{forest} nominal tree
[NP
 [{\synnode{Head}{Nom}}, head edge
  [{\synnode{Head}{DP}}, head edge
   [{\synnode{Head}{D}}, head edge [\mention{some}]]]]]
\end{forest}
\end{minipage}
\caption{\mention{Some apples} (a), and \mention{some} standing alone under fusion (b) and ordinary Head (c). Function labels sit above category labels, and thick lines join heads to their parents. In (b) the DP has two parents, with a line from the NP for its determiner function and one from the nominal for its head function. In (c) it's the head of the nominal and nothing else. Each tree shows the DP even where it contains only D, so that the same constituent appears in all three.}\label{fig:analyses}
\end{figure}

Without a noun, as in \mention{I'll take \bk{\fu{some}}}, the nominal still needs a head. \textit{CGEL}'s answer is \term{fusion of functions}: the DP functions as determiner and as head at once, a fused determiner--head (Det--Head; Figure~\ref{fig:analyses}b) (pp.~410--421). After a determiner, as in \mention{the lucky few}, the word is a fused modifier--head (Mod--Head). PHP make the device explicit and list where English uses it: \enquote{The subtypes of dependent involved are determiner, modifier and, in fused relative clauses, subject or prenucleus} (p.~571; the prenucleus is the position of a fronted relative phrase). Sections~\ref{sec:det-head}, \ref{sec:mod-head} and~\ref{sec:fused-relatives} take them in that order, and Section~\ref{sec:once} takes PHP's own case, \mention{once}, \mention{twice}, and \mention{thrice}.

Fusion, as PHP define it, departs from canonical structure in two ways. First, one constituent has two functions, so that it \enquote{will in effect have two mothers}, which they call \enquote{the major departure from canonical principles of structure} (p.~570). The structure is then a directed acyclic graph rather than a tree \citep{pullumRogers2008Expressive}. In Det--Head the shared DP has two parents, the NP and the nominal; in Mod--Head both functions belong to one nominal (pp.~582, 588). Second, an NP can have an ultimate head of another category: the fused node is \enquote{a point in the normal path of headedness at which there can be a change of category} (p.~569).

\term{Ordinary Head} keeps the second departure and drops the first. Whatever phrase \textit{CGEL} fuses functions as head of the nominal and nothing else: a DP for \mention{some} (Figure~\ref{fig:analyses}c) and for \mention{few} in \mention{the lucky few}, an AdjP in \mention{the rich}. For adjectives this is close to the traditional analysis, though there the adjective heads the noun phrase itself: \enquote{Adjectives can function as heads of noun phrases} \citep[421]{quirk1985}. A determiner before it, as \mention{the} in \mention{the rich}, is the NP's determiner as usual, and the nominal keeps its other dependents: in \mention{few of them} the \mention{of}-phrase is a complement in the nominal headed by \mention{few}, as in \textit{CGEL}'s own partitive tree (pp.~411--412). What fusion states as a licence for a fused function, ordinary Head states as a \term{permission}: a word may head a nominal, with or without a determiner before it. \mention{Some} may head one with no determiner, and \mention{rich} one after \mention{the}.

\textit{CGEL} has a term for a head of this kind: the independent genitive pronouns \enquote{occur in fused heads and also pure heads} (p.~410), a pure head being one with no second function. It also considers the analysis, as \enquote{change of function without change of category}, with \mention{many} a determiner in \mention{\bk{Many} people would agree} and a head in \mention{\bk{\fu{Many}} would agree} (p.~421).

The change of category needs no new kind of structure. PHP say it \enquote{will be independently required in the case of hybrid constructions} (p.~571), and their tree for \mention{there's \bk{no stopping him}} has a VP as head of an NP whose determiner is \mention{no}: \enquote{an unexpected category appears as head \ldots\ However, there is no fusion of functions in this case} (p.~600). The tree shares the two properties that matter here, an unexpected category as head and no constituent with two functions, though the VP heads the NP directly, with no nominal between. PHP find hybrids \enquote{only marginally in English} (p.~566). But ordinary Head changes category only where fusion does, and at the same node, so it uses no departure that fusion doesn't also use, and uses it no more widely.

Both analyses keep three commitments that \textit{CGEL} argues for, and the comparison holds them fixed. First, in \mention{some apples} the head is \mention{apples}, not \mention{some} (Figure~\ref{fig:analyses}a). The noun is the NP's ultimate head, and the determinative heads a DP that functions only as determiner. Second, the determinative keeps its category when it stands alone: it takes the dependents it takes before a noun (\mention{not much cheese}, \mention{not much of the cheese}), which \enquote{argues against any part-of-speech distinction between the two uses} (\textit{CGEL}, p.~420). Third, there's no silent head noun. Analyses that give up one of these avoid fusion by other means: a determiner phrase headed by the determinative, or a silent noun.\footnote{On the DP hypothesis the determinative heads the phrase \citep[see][]{millerpullum2022tete,salzmann2020npvsdp}, and pronouns can be analysed as a kind of definite article \citep{postal1966}. A silent noun is recovered from an antecedent on ellipsis analyses \citep{lobeck1995ellipsis} and needs none on empty-noun analyses \citep{kester1996adjectival,panagiotidis2003empty,saab2019nominal,kayne2007several,aarts2024fused}. \textcite{wu2021exocentric} puts an empty noun after \mention{the rich} and after a determinative standing alone, rejecting adjective and determinative heads because a head should share its phrase's category and a bare adjective can't show number (pp.~19--20, 118, 173). Ordinary Head makes the adjective head of the nominal, not of the NP, and \mention{the rich} is plural with no plural ending on every analysis, as Wu grants (p.~118 n.~14). \textit{CGEL} gives its grounds against a silent head on pp.~420--421.} This paper doesn't take those analyses on. Its result is conditional: given \textit{CGEL}'s commitments, the balance favours fusion only in fused relatives. Holding \textit{CGEL}'s categories fixed isolates the question of functions; it doesn't endorse them. Indeed, \textcite{reynolds2026determinativesnouns} argues that determinatives are a kind of noun, beside common nouns, proper nouns, and pronouns. On that classification a determinative heading a nominal is a noun heading it, so the second departure disappears for determinatives on both analyses, leaving ordinary Head with no departure for them and fusion with the first.

Since both analyses describe the same strings, the comparison needs a criterion. Interpretation can't supply it, since the two assign the same readings, one by a second function, the other by what the context supplies. What's left is form, which is what \textit{CGEL} rests the determiner function on when it argues from the genitive (p.~420). PHP call a constituent with two functions a departure from canonical structure, so I take the burden to lie with the second function. Where the same category and head relation account for an expression, a property of form counts as evidence for a second function if it varies with that function and not with the word alone. By form I mean what can be observed of an expression: the shapes of its words (a case such as \mention{whom}, or a variant such as \mention{none}, which occurs only without a noun after it), their order, and which dependents and determiners occur with each. Interpretation alone isn't form. \textit{CGEL} keeps the two apart elsewhere. In \mention{She considered Ed \bk{quite competent}}, the adjective phrase is predicated of \mention{Ed}, which is the object of \mention{considered}, and \textit{CGEL} calls \mention{Ed} the adjective phrase's \term{predicand}, a semantic relation, not its subject, a syntactic one (pp.~217, 252).

The criterion answers \textit{CGEL}'s objection to ordinary Head. \mention{Nylon} in \mention{I prefer cotton shirts to \bk{\fu{nylon}}} is one of \textit{CGEL}'s fused modifier--heads, of the type denoting composition (p.~417). \textit{CGEL} rejects analysing it as a noun functioning simply as head, since that \enquote{would not distinguish this construction from \mention{I prefer cotton to nylon}}, where \mention{nylon} isn't \enquote{understood as having modifier function} (p.~421). The difference is real. But the word is the same in both sentences, and so is its form: bare \mention{nylon} as the object of \mention{to}. What differs is what the context supplies.

The one formal property \textit{CGEL} ties to these modifier--heads is their determiners. The NPs of \mention{nylon}'s type \enquote{are either bare or have a definite determiner}, with \enquote{significant loss of acceptability when the determiner is indefinite} (p.~417), and nouns of other types are excluded, as in \ungram{\mention{Bill likes the linguistics lecturer, but I prefer \bk{the \fu{sociology}}}} (p.~415). The examples in (\ref{ex:dog}) cross both lines.

\ea\label{ex:dog}
\ea \mention{You take the cat shirt and I'll take \bk{the \fu{dog}}.}
\ex \ungram{\mention{You take the cat shirt and I'll take \bk{\fu{dog}}.}}
\ex \mention{You take the cotton shirt and I'll take \bk{\fu{nylon}}.}
\ex \mention{a misdemeanor sentence rather than \bk{a \fu{felony}}} (\textit{Arizona Daily Star}, 2019, COCA)
\z\z

\mention{Dog} and \mention{felony} belong to no type \textit{CGEL} lists, and they work. Bare \mention{dog} is out and \mention{a felony} is fine, as for a singular count noun generally. Non-count \mention{nylon} stands bare for a single item and resists \mention{a} (\marg{\mention{I bought a cotton shirt rather than \bk{a \fu{nylon}}}}, p.~417). So the determiners follow the noun's own countability. Both analyses make the noun the head, so both predict that. Neither kind of silent head does: every antecedent in (\ref{ex:dog}) is a count noun, so a head that took its features would reject bare \mention{nylon}, and a head with none of its own wouldn't exclude bare \mention{dog}. What the facts don't show is anything that varies with the modifier function, and the restriction \textit{CGEL} states for it doesn't hold.\footnote{Dependents give the same result. \textcite[28--29]{matthews2014} notes that material nouns used as modifiers take \mention{partly} and \mention{largely} (\mention{a partly cotton shirt}) and suggests they're adjectives there. But uses with no modifier function take them too (\mention{I prefer cotton to \bk{partly nylon}}; \mention{This lava is \bk{partly magma}}, conserve-energy-future.com), so \mention{partly} isn't tied to the modifier function.} On this criterion, nothing here counts as evidence for the modifier function.\footnote{\textit{CGEL}'s other example, \mention{Mary earns \bk{\fu{double}}} (p.~421), fuses a noun with a predeterminer modifier (p.~418), and \mention{double} gets the treatment \mention{nylon} gets. Its category fits. Before a definite determiner, \textit{CGEL}'s predeterminers are nouns and NPs (\mention{half the cake}, \mention{double the number of applications}) or the adverb \mention{twice} (p.~434), and my searches of the Corpus of Contemporary American English (COCA; \citealp{davies2008cocaonline}) found no adjective phrase there, apart from one marginal \mention{400-fold the amount} (a blog).}

\textit{CGEL} also names a form that a second function would show: a fused determiner--head should resist the genitive, because \enquote{for full acceptability its case should match both functions} (p.~420), as the form of \mention{whoever} or \mention{whomever} reflects its function inside a fused relative (pp.~1073--1074). Sections~\ref{sec:det-head} and~\ref{sec:fused-relatives} run that test.

The criterion says what counts as evidence for a second function, not how much is enough. A fact counts if it varies with the function and the word alone doesn't explain it. Whether it outweighs the extra structure is a separate question, and a fact that something else could also explain weighs less. So the later sections compare the analyses construction by construction: what each predicts, where the data go against a prediction, what each has to state instead, and what the extra structure buys. Where the data can't separate the analyses, I say so, since finding no evidence for a second function isn't finding evidence against it. The standard also applies to what this paper adds. Its analysis of \mention{a few} rests on number and on the structure of numerals (Section~\ref{sec:mod-head}), and in fused relatives the permissions a head analysis needs count against it (Section~\ref{sec:fused-relatives}).

A note on data. Examples are \textit{CGEL}'s or PHP's where cited, attested where a source is given, and otherwise constructed. Judgments not attributed to a source are mine. Corpus lines come from COCA, the Corpus of News on the Web (NOW; \citealp{davies2016now}), and the iWeb corpus \citep{davies2018iweb}, searched by hand at english-corpora.org in September and October 2026; where a display was capped, counts are of the lines shown. Some sets of lines were classified by a large language model (Claude Opus~5.5), and each line's label is kept so that the classification can be checked. The genitive lines in Section~\ref{sec:det-head} were labelled twice, independently, with the same result. I checked the labels behind Sections~\ref{sec:det-head} and~\ref{sec:fused-relatives} by hand (for the earlier searches, a sample of 105) and changed five, none affecting a figure the paper reports. The record behind each corpus claim, with its query, date, counts, and labels (but not the corpus lines), is listed in an \href{https://github.com/BrettRey/fusion-of-functions-revisited/blob/master/EVIDENCE.md}{evidence index}.
